\section{Discussion}

\subsection{Why the Hybrid Signal Works}
The two reward components are \emph{complementary}: NLI entailment $P_\text{NLI}(E \models A)$ verifies that the explanation logically implies the correct option, while the verifier score $S_\text{Verifier}$ captures clarity, structure, and pedagogical fit that NLI cannot detect. Used alone, each degenerates---pure NLI converges to short answer-repeating snippets, pure verifier reintroduces the verbosity bias---so we score candidates by their conjunction (Eq.~\ref{eq:hybrid_reward}). This joint constraint filters both degenerate extremes and steers the policy toward the upper-right quadrant of the alignment frontier (Figure~\ref{fig:alignment_tax}). The verifier-blindness finding (Section~4.2: chosen-NLI / rejected-NLI gap of $17$--$18\times$ alongside indistinguishable chosen/rejected verifier scores) replicates across LLaMA-2-13B, Qwen3-8B, Gemma 4 E4B-it, and the Tier-B ablation, indicating that NLI is the discriminating signal and that the verifier acts as a fluency regulariser rather than a logical scorer.

\subsection{Architecture Effects and Entailment Headroom}
\label{sec:arch-effects}
Across the three architectures, the marginal NLI gain from Hybrid-DPO is inversely related to the SFT NLI baseline. LLaMA-2-13B SFT scores $0.05$--$0.09$ on NLI and RLearner-LLM lifts it to $0.32$--$0.43$ ($4$--$6\times$); Qwen3-8B SFT is already at $0.16$--$0.32$ and gains shift onto ACR ($+9$ to $+29$\,pp); Gemma 4 E4B-it (a $4.5$B-effective-parameter dense model with Per-Layer Embeddings) improves NLI on $4$ of $5$ domains, including the largest single-domain lift in this paper ($0.1604 \to 0.3892$ on UK Medicine Year 2, virtually matching the LLaMA-2-13B peak with substantially fewer effective parameters), and is the first single-pass RL result in our study to surpass the iterative ILearner-LLM (K=5) benchmark on Auckland Law. The one Gemma 4 non-improvement is Sydney Biology, where the SFT NLI baseline is already the highest of our three architectures ($0.2469$). Together these patterns suggest that Hybrid-DPO's effectiveness tracks \emph{entailment headroom}---the gap between the SFT policy's current NLI and its attainable ceiling on a given domain---rather than raw model scale.

On four of five domains the Gemma 4 RLearner's NLI gain comes with a modest ACR reduction ($-10.3$ to $-1.3$\,pp): the ACR gate enforces $\mathrm{ACR}(E)\geq 0.5$ on every chosen pair, but the conjunctive reward then optimises NLI \emph{subject to} that floor rather than pushing past it. A stronger ACR weight (or per-domain preference data) is a natural extension to recover Qwen3's ACR profile without sacrificing Gemma 4's NLI gains.

\subsection{Limitations}
\label{sec:discussion-limits}
Three limitations bound our results. (i) On Auckland Law (LLaMA-2), single-pass DPO reaches NLI $0.3229$ but is surpassed by the iterative ILearner-LLM (K=5) at $0.3996$ (Gemma 4 closes this gap; \S\ref{sec:arch-effects}), suggesting a hybrid regime that adds targeted iteration in highly structured domains as future work. (ii) Our NLI scorer is the compact \texttt{cross-encoder/nli-deberta-v3-small}; a stronger classifier (e.g.\ DeBERTa-v3-large) would likely yield more discriminative preference labels at higher compute. (iii) The $13{,}211$-example SFT corpus is authored by undergraduate students on PeerWise, not subject-matter experts; partial justifications and occasional factual inaccuracies place a soft upper bound on overlap-style metrics---this provenance is itself the motivation for the Hybrid reward, which grounds the preference signal in NLI against the correct-option text rather than in surface similarity. Re-running on an expert-curated corpus is a clear next step.

\subsubsection*{Robustness to Question-Quality Filtering}

To check whether the gains in Section~4 depend on the breadth of the question pool, we re-ran the entire Cardiff Biology pipeline on a stricter ``Tier-B'' subset that adds \texttt{deleted}\,${=}$\,0 and a $5$-star endorsement share $\geq 0.10$ to the default filter, shrinking the Cardiff training pool $7\times$ ($7{,}309 \to 1{,}041$) with all other hyperparameters held fixed (Table~\ref{tab:tierB_ablation}).

\begin{table}[h]
\centering
\small
\resizebox{\textwidth}{!}{
\begin{tabular}{l|c|c|c|c|c|c|c}
\toprule
Cardiff filter & $N$ & SFT NLI & DPO BLEU & DPO BERT(Ans) & DPO ACR & DPO NLI & DPO time(s) \\
\midrule
Default (Section~4) & $7{,}309$ & 0.2117 & 0.0303 & 0.8426 & 0.6497 & 0.3505 & \textbf{4.76} \\
\textbf{Tier-B (stricter)} & $1{,}041$ & \textbf{0.2059} & \textbf{0.0381} & \textbf{0.8696} & \textbf{0.7823} & \textbf{0.5202} & 7.62 \\
\midrule
$\Delta$ (Tier-B $-$ default) & $-86\%$ & $-0.006$ & $+0.008$ & $+0.027$ & $+0.133$ & \textbf{+0.170 ($+48\%$)} & $+2.9$s \\
\bottomrule
\end{tabular}
}
\caption{Cardiff Biology robustness to question-quality filtering. Tier-B adds \texttt{deleted}=0 and a $5$-star endorsement share $\geq 0.10$ to the default filter, shrinking the Cardiff training corpus by $7.0\times$. The SFT baseline is essentially unchanged (NLI $0.2117 \to 0.2059$), but Hybrid-DPO trained on the smaller Tier-B corpus reaches \emph{higher} NLI ($0.3505 \to 0.5202$), ACR ($0.6497 \to 0.7823$), and answer-anchored BERTScore than under the default filter, with comparable BLEU and BERT(Stu).}
\label{tab:tierB_ablation}
\end{table}

Despite the $7\times$ smaller training pool, Tier-B Hybrid-DPO \emph{exceeds} the default-filter result on every quality metric, lifting NLI by $+48\%$ ($0.3505 \to 0.5202$) and ACR by $+13.3$\,pp. We attribute this to a cleaner upstream signal: the Tier-B chosen/rejected NLI gap is $18\times$ (vs $17\times$ at default) while verifier scores remain indistinguishable ($3.24$ vs $3.27$)---community endorsement filters out pairs whose chosen explanation comes from ambiguous or low-consensus questions, sharpening DPO's gradient. A strict-discriminator filter (dropping questions whose author-marked answer disagrees with the modal student answer) needs per-option response counts that are not in our metadata; given the corpus-size-insensitive Tier-B result, we expect it to extend rather than reverse this trend.

\subsection{Implications and Verbosity-Bias Verdict}
Two practitioner-facing implications follow. First, in any knowledge-intensive generation task where outputs have a verifiable logical property (medical diagnosis, mathematical proof, code correctness, legal reasoning), the standard practice of using stylistic human preferences as the DPO reward leaves a substantial logical alignment gap; \emph{the reward-signal specification is the most critical design decision} in applying DPO to such tasks. Second, the $69\%$ LLM-judge win rate for the verbose LLaMA-2 SFT baseline (despite its NLI of $0.055$) is not a calibration glitch but a systematic preference of LLM-as-a-judge for hedging/enumerative/formal register over logical content. Our $95\%$ win rate of RLearner-LLM (Qwen3) against GPT-4o-mini in blind pairwise evaluation shows the bias is overcomeable, but only once the policy has been aligned to produce concise, logically dense explanations that defeat the verbosity heuristic. This argues for logic-aware automatic metrics (NLI, ACR) as primary evaluation criteria for educational NLG, rather than relying solely on LLM judges or surface overlap metrics.
